% concatenated from arXiv_locatedness,_convexity,_and_measurability_230726.tex
%2multibyte Version: 5.50.0.2960 CodePage: 1253
\documentclass[11pt]{article}%
\usepackage{amssymb}
\usepackage{amsmath}
\usepackage{euler}
\usepackage{amsfonts}
\usepackage{graphicx}%
\setcounter{MaxMatrixCols}{30}
%TCIDATA{OutputFilter=latex2.dll}
%TCIDATA{Version=5.50.0.2960}
%TCIDATA{Codepage=1253}
%TCIDATA{CSTFile=40 LaTeX article.cst}
%TCIDATA{Created=Monday, July 13, 2026 11:40:47}
%TCIDATA{LastRevised=Thursday, July 23, 2026 10:45:01}
%TCIDATA{<META NAME="GraphicsSave" CONTENT="32">}
%TCIDATA{<META NAME="SaveForMode" CONTENT="1">}
%TCIDATA{BibliographyScheme=Manual}
%TCIDATA{<META NAME="DocumentShell" CONTENT="Standard LaTeX\Blank - Standard LaTeX Article">}
%TCIDATA{Language=American English}
%BeginMSIPreambleData
\providecommand{\U}[1]{\protect\rule{.1in}{.1in}}
%EndMSIPreambleData
\newtheorem{theorem}{Theorem}
\newtheorem{acknowledgement}[theorem]{Acknowledgement}
\newtheorem{algorithm}[theorem]{Algorithm}
\newtheorem{axiom}[theorem]{Axiom}
\newtheorem{case}[theorem]{Case}
\newtheorem{claim}[theorem]{Claim}
\newtheorem{conclusion}[theorem]{Conclusion}
\newtheorem{condition}[theorem]{Condition}
\newtheorem{conjecture}[theorem]{Conjecture}
\newtheorem{corollary}[theorem]{Corollary}
\newtheorem{criterion}[theorem]{Criterion}
\newtheorem{definition}[theorem]{Definition}
\newtheorem{example}[theorem]{Example}
\newtheorem{exercise}[theorem]{Exercise}
\newtheorem{lemma}[theorem]{Lemma}
\newtheorem{notation}[theorem]{Notation}
\newtheorem{problem}[theorem]{Problem}
\newtheorem{proposition}[theorem]{Proposition}
\newtheorem{remark}[theorem]{Remark}
\newtheorem{solution}[theorem]{Solution}
\newtheorem{summary}[theorem]{Summary}
\newenvironment{proof}[1][Proof]{\noindent\textbf{#1.} }{\ \rule{0.5em}{0.5em}}
\begin{document}

\date{}
\title{\textsf{Locatedness, Convexity, and Measurability in }$\mathbb{R}^{N}$}
\author{\textsf{Douglas S. Bridges}}
\maketitle

\begin{abstract}%
%TCIMACRO{\TeXButton{noindent}{\noindent}}%
%BeginExpansion
\noindent
%EndExpansion
\textsf{We prove, within Bishop's constructive framework \cite{Bishop},
stronger forms of two theorems in \cite{Br88} about the locatedness of certain
Lebesgue measurable complemented sets in }$R^{N}$\textsf{. In doing so, we
expand and improve the proofs of several preliminary results on Lebesgue
measurability and convexity that are need en route to our main theorems.}

\end{abstract}

%

%TCIMACRO{\TeXButton{depth 1}{\setcounter{secnumdepth}{0}}}%
%BeginExpansion
\setcounter{secnumdepth}{0}%
%EndExpansion%
%TCIMACRO{\TeXButton{sf}{\normalfont\sf}}%
%BeginExpansion
\normalfont\sf
%EndExpansion
%

%TCIMACRO{\TeXButton{bigskip}{\bigskip}}%
%BeginExpansion
\bigskip
%EndExpansion
%

%TCIMACRO{\TeXButton{noindent}{\noindent}}%
%BeginExpansion
\noindent
%EndExpansion
The framework of this paper\footnote{%
%TCIMACRO{\TeXButton{sf}{\normalfont\sf}}%
%BeginExpansion
\normalfont\sf
%EndExpansion
\textbf{Keywords: \ }measurable, convex, located \ \textbf{MSC
classifications: \ }26E40, 03F60} is Bishop's constructive mathematics. We
assume some familiarity with, or at least access to, relevant parts of the
constructive integration theory found in Chapter 6 of \cite{BB}. Additional
sources for constructive analysis are \cite{Bishop}, \cite{BVtech}, and---the
most up-to-date reference---Ishihara's excellent Chapters 8 and 9 in
\cite{Handbook}.

Throughout the paper, $N$ will be a fixed positive integer, $\rho$ the
Euclidean metric on $\mathbb{R}^{N}$, and $\mu$ Lebesgue measure on
$\mathbb{R}^{N}$. The words\emph{ integrable} and \emph{measurable} will refer
exclusively to Lebesgue measure theory.

A \textsc{complemented set} in $\mathbb{R}^{N}$ has the form $\mathbf{A}%
\equiv\left(  A^{1},A^{0}\right)  $ of subsets of $\mathbb{R}^{N}$ such
that\footnote{%
%TCIMACRO{\TeXButton{sf}{\normalfont\sf} }%
%BeginExpansion
\normalfont\sf
%EndExpansion
For real numbers $x,y$ we have $x\neq y$ if and only if $\left\vert
x-y\right\vert >0$. Constructively, $\neq$ is equivalent to the negation of
equality if and only if we accept Markov's principle, which is widely regarded
as of at least dubious constructive validity \cite[pp. 10--11]{BVtech}.}
$x\neq y$ whenever $x\in A^{1}$ and $y\in A^{0}$. The characteristic function
of $\mathbf{A}$ is then the mapping $\chi_{\mathbf{A}}$ with domain $A^{1}\cup
A^{0}$ such that $\chi_{\mathbf{A}}(x)=k$ if $x\in A^{k}$. We say that
$\mathbf{A}$ is integrable/measurable if $\chi_{A}$ is integrable/measurable,
in which case its domain is a full subset of $\mathbb{R}^{N}$. We say that
$\mathbf{A}$ \textsc{has positive (possibly infinite) measure} if it is
measurable and there exists an integrable complemented set $\mathbf{B}$ in
$\mathbb{R}^{N}$ such that $\mathbf{B}<\mathbf{A}$ (that is, $\chi
_{\mathbf{B}}\leqslant\chi_{\mathbf{A}}$) and $\mu(\mathbf{B})>0$.

Our objective is to present stronger versions of \cite[Thms 2 and 4]{Br88}---namely,

\begin{theorem}
\label{jul21t1}Let $\mathbf{S}$ be a measurable complemented set in
$\mathbb{R}^{N}$ with positive (possibly infinite) measure, such that $S^{1}$
is convex. Then $S^{1}$ is located in $\mathbb{R}^{N}$.
\end{theorem}

\begin{theorem}
\label{jul21t2}Let \textbf{$S$} be a measurable complemented set in
$\mathbb{R}^{N}$\emph{ }such that $S^{1}$ is convex and $%
%TCIMACRO{\TeXButton{mathsim}{\mathord{\sim}}}%
%BeginExpansion
\mathord{\sim}%
%EndExpansion
S^{1}\ $is inhabited. Then $%
%TCIMACRO{\TeXButton{mathsim}{\mathord{\sim}}}%
%BeginExpansion
\mathord{\sim}%
%EndExpansion
S^{1}$ is located in $\mathbb{R}^{N}$.
\end{theorem}

%

%TCIMACRO{\TeXButton{noindent}{\noindent}}%
%BeginExpansion
\noindent
%EndExpansion
Here,%
\[%
%TCIMACRO{\TeXButton{mathsim}{\mathord{\sim}}}%
%BeginExpansion
\mathord{\sim}%
%EndExpansion
S^{1}\equiv\left\{  x\in\mathbb{R}^{N}:x\neq s\text{ for all }s\in
S^{1}\right\}  ,
\]
the \textsc{complement} of $S^{1}$ in $\mathbb{R}^{N}$, not to be confused
with the \textsc{metric complement }of $S^{1}$ in $\mathbb{R}^{N}$,%
\[
-S^{1}\equiv\left\{  x\in\mathbb{R}^{N}:\exists_{r>0}\forall_{s\in S^{1}%
}\left(  \left\vert x-s\right\vert \geq r\right)  \right\}  .
\]
%

%TCIMACRO{\TeXButton{medskip}{\medskip}}%
%BeginExpansion
\medskip
%EndExpansion


We begin with some measure theory. By an $n$\textsc{-dimensional
rectangle}\footnote{%
%TCIMACRO{\TeXButton{sf}{\normalfont\sf}}%
%BeginExpansion
\normalfont\sf
%EndExpansion
https://grokipedia.com/page/Hyperrectangle} we mean the Cartesian product
$S\equiv I_{1}\times\cdots\times I_{n}$ of $n$ compact intervals in
$\mathbb{R}$. It readily follows from Fubini's theorem \cite[Ch. 6, (9.7)]{BB}
that the complemented set $\left(  S,-S\right)  $ is integrable, with measure
$\left\vert I_{1}\right\vert \times\cdots\times\left\vert I_{n}\right\vert $,
where $\left\vert I\right\vert $ is the length of a compact interval in
$\mathbb{R}$. This measure is the hypervolume of the rectangle $S$.

\begin{lemma}
\label{jul03l1}Let $V$ be an $n$-dimensional subspace of $\mathbb{R}^{N}$,
where $1\leqslant n<N$. Then $\mathbf{V}\equiv\left(  V,-V\right)  $ is an
integrable complemented subset of $\mathbb{R}^{N}$, and $\mu(\mathbf{V})=0$.
\end{lemma}

\begin{proof}
Let $\left\{  e_{1},\ldots,e_{N}\right\}  $ be the standard orthonormal basis
of $\mathbb{R}^{N}$. We may assume that $V=\mathsf{span}\left\{  e_{1}%
,\ldots,e_{n}\right\}  $. For each positive integer $k$,%
\[
A_{k}\equiv\underset{n~\text{\textsf{terms}}}{\underbrace{\left[  -k,k\right]
\times\cdots\times\left[  -k,k\right]  }}\times\underset{N-n\text{
\textsf{terms}}}{\underbrace{\left\{  0\right\}  \times\cdots\times\{0\}}}%
\]
is a rectangle in $\mathbb{R}^{N}$, so, referring to the definition above, we
see that the complemented set%
\[
\mathbf{A}_{k}\equiv\left(  A_{k},-A_{k}\right)
\]
is integrable, with measure $0$. Since $\sum_{k=1}^{\infty}\mu(\mathbf{A}%
_{k})=0$, it follows from \cite[Ch. 6, (3.10)]{BB} that $%
%TCIMACRO{\dbigvee \nolimits_{k\geqslant1}}%
%BeginExpansion
{\displaystyle\bigvee\nolimits_{k\geqslant1}}
%EndExpansion
\mathbf{A}_{k}$ is integrable and $\mu\left(
%TCIMACRO{\dbigvee \nolimits_{k\geqslant1}}%
%BeginExpansion
{\displaystyle\bigvee\nolimits_{k\geqslant1}}
%EndExpansion
\mathbf{A}_{k}\right)  =0$. It is routine to verify that $%
%TCIMACRO{\dbigvee \nolimits_{k\geqslant1}}%
%BeginExpansion
{\displaystyle\bigvee\nolimits_{k\geqslant1}}
%EndExpansion
\mathbf{A}_{k}=(V,-V)$.
\end{proof}

\begin{lemma}
\label{jul03l2}Let $\mathbf{A}$ be an integrable complemented set in
$\mathbb{R}^{N}$. Let $V$ be an $n$-dimensional subspace of $\mathbb{R}^{N}$,
where $1\leqslant n<N$, and let $\mathbf{V}$ be the complemented set $\left(
V,-V\right)  $. Then $\mathbf{V}\wedge\mathbf{A}$ is integrable, with measure
$0$.
\end{lemma}

\begin{proof}
By Lemma \ref{jul03l1}, $\mathbf{V}$ is integrable and $\mu(\mathbf{V})=0$; so
$\mathbf{V}\wedge\mathbf{A}$ is integrable, by \cite[Ch.6, (3.5)]{BB};
$\mu\left(  \mathbf{V}\wedge\mathbf{A}\right)  \leqslant\mu(\mathbf{V})$, by
\cite[Ch. 6, (3.7)]{BB}; and therefore $\mu\left(  \mathbf{V}\wedge
\mathbf{A}\right)  =0$.
\end{proof}

\begin{lemma}
\label{jul21l1}Let $\mathbf{A}$ and $\mathbf{B}$ be, respectively, a
measurable and an integrable complemented set in $\mathbb{R}^{N}$. Then
$\mathbf{A\wedge B}$ is integrable and $\mu(\mathbf{A\wedge B})\leqslant
\mu(\mathbf{B})$.
\end{lemma}

\begin{proof}
Since $\mathbf{B}$ is also measurable, this is a simple consequence of
\cite[Ch. 6, (7.7) and (7.11)]{BB}.
\end{proof}

%

%TCIMACRO{\TeXButton{medskip}{\medskip}}%
%BeginExpansion
\medskip
%EndExpansion


For each $x\in\mathbb{R}^{N}$ and each $r>0$, denote by $B(x,r)$ the open ball
in $\mathbb{R}^{N}$ with centre $x$ and radius $r$, and define the
complemented set%
\[
\mathbf{B}(x,r)\equiv\left(  B(x,r),-B(x,r)\right)  .
\]
Then $\mathbf{B}(x,r)$ is integrable, $\mu(\mathbf{B}(x,r))$ being the
hypervolume of a ball of radius $r$ in $\mathbb{R}^{N}$.

\begin{lemma}
\label{jul22l1}Let $\mathbf{S}$ be a measurable complemented set in
$\mathbb{R}^{N}$ such that $S^{1}$ is bounded. Then $\mathbf{S}$ is integrable.
\end{lemma}

\begin{proof}
There exists $r>0$ such that $S^{1}\subset B(0,r)$. Then $\mathbf{S}%
=\mathbf{S\wedge B}(0,r)$ and is therefore integrable, by Lemma \ref{jul21l1}.
\end{proof}

\begin{lemma}
\label{jul02l1}Let $\mathbf{A}$ be a measurable complemented set in
$\mathbb{R}^{N}$ with positive (possibly infinite) measure. Then $A^{1}$
contains a basis of $\mathbb{R}^{N}$.
\end{lemma}

\begin{proof}
Suppose first that $\mathbf{A}$ is integrable, so that $\mu(\mathbf{A})>0$. By
\cite[Ch. 6, (3.4)]{BB}, $A^{1}$ is inhabited by some point $a$. Since
$\left(  \left\{  a\right\}  ,-\left\{  a\right\}  \right)  $ is integrable,
with measure $0$, we see from \cite[Ch. 6, (3.6)]{BB} that $\mathbf{A}-\left(
\left\{  a\right\}  ,-\left\{  a\right\}  \right)  $ is integrable and has
measure%
\[
\mu(\mathbf{A})-\mu\left(  \left(  \left\{  a\right\}  ,-\left\{  a\right\}
\right)  \right)  =\mu(\mathbf{A})>0\text{.}%
\]
Thus, by \cite[Ch. 6, (3.4)]{BB}, $A^{1}-\left\{  a\right\}  $ is inhabited by
some point $b$. Since $a\neq b$, it follows that $A^{1}$ contains a nonzero
point $e_{1}$ of $\mathbb{R}^{N}$. Now assume that for some $n<N$, we have
constructed linearly independent vectors $e_{1},\ldots,e_{n}$ belonging to
$A^{1}$. Let $V$ be the $n$-dimensional subspace of $\mathbb{R}^{N}$ spanned
by these vectors, and let $\mathbf{V}=\left(  V,-V\right)  $. By Lemma
\ref{jul03l2}, $\mathbf{V}\wedge\mathbf{A}$ is $\mu$-integrable and
$\mu(\mathbf{V}\wedge\mathbf{A})=0$. By \cite[Ch. 6, (3.6)]{BB},
$\mathbf{A}-\mathbf{V}\wedge\mathbf{A}$ is $\mu$-integrable, and%
\[
\mu(\mathbf{A}-\mathbf{V}\wedge\mathbf{A})=\mu(\mathbf{A})-\mu(\mathbf{V}%
\wedge\mathbf{A})>0.
\]
Hence, again by \cite[Ch. 6, (3.6)]{BB}, there exists $e_{n+1}\in A^{1}\cap
-V$. Then $\rho(e_{n+1},V)>0\,$, so by Lemma 4.1.10 of \cite{BVtech}, the
vectors $e_{1},\ldots,e_{n+1}$ in $A^{1}$ are linearly independent. The
desired conclusion for integrable $\mathbf{A}$ now follows by induction. In
the general case, $\mathbf{A}$ has an integrable complemented subset
$\mathbf{B}$ with $\mu(\mathbf{B})>0$, so by the foregoing, $B^{1}$, and
therefore $A^{1}$, contains a basis of $\mathbb{R}^{N}$.
\end{proof}

%

%TCIMACRO{\TeXButton{medskip}{\medskip}}%
%BeginExpansion
\medskip
%EndExpansion


Turning now to convexity, for convenience we state here Lemma 5.1.1 of
\cite{BVtech}.

\begin{lemma}
\label{jul14l1}Let $C$ be a convex subset of a normed space $X$, let $\xi\in
C^{\circ}$, and let $r>0$ be such that $B(\xi,r)\subset C.$ Let $x\neq\xi,$
$0<t<1,$ and $y=t\xi+\left(  1-t\right)  x.$ If the ball $B(x,tr)$ intersects
$C,$ then $B(y,t^{2}r)\subset C$.
\end{lemma}

%

%TCIMACRO{\TeXButton{medskip}{\medskip}}%
%BeginExpansion
\medskip
%EndExpansion


Let $S$ be a subset of a metric space $X$. We say that $S^{\circ}$
is\textbf{\ }\textsc{uniformly dense in} $S$ if for each $\varepsilon>0$ there
exists $\delta>0$ such that if $x\in S$, then $S\cap B(x,\varepsilon)$
contains a ball in $X$ of radius $\delta$. In that case, $S^{\circ}$ is
clearly dense in $S$.

The proofs of the next two lemmas improve upon those of \cite[Lemmas 9 and
2]{Br88}.

\begin{lemma}
\label{0105a2}Let $C$ be a bounded convex set with inhabited interior in a
normed space. Then $C^{\circ}$ is uniformly dense in $C$ \emph{\cite[Lemma
9]{Br88}}.
\end{lemma}

\begin{proof}
\label{ere}Fix\textsf{\ }$\xi\in C$ and $r>0$ such that $B(\xi,r)\subset C$.
There exists $R>r$ such that $C\subset B(\xi,R)$. Given $\varepsilon$ with
$0<\varepsilon<R$, let $t=\varepsilon/2R$ and $\delta=t^{2}r$; then $0<t<1$.
For each $x\in C$ either $\left\Vert x-\xi\right\Vert <r/4$ or $x\neq\xi$. In
the first case, for each $y\in B(x,\delta)$ we have%
\[
\left\Vert y-\xi\right\Vert \leqslant\left\Vert y-x\right\Vert +\left\Vert
x-\xi\right\Vert <t^{2}r+\frac{r}{4}<\left(  \frac{\varepsilon}{R}\right)
^{2}\frac{r}{4}+\frac{r}{4}<r,
\]
so $B(x,\delta)\subset B(\xi,r)\subset C$. In the case $x\neq\xi$, write
$z\equiv t\xi+(1-t)x$. Since $x\in C\cap B(x,tr)$, Lemma \ref{jul14l1} tells
us that $B(z,\delta)\subset C$. On the other hand, if $y\in B(z,\delta)$, then%
\[
\left\Vert x-y\right\Vert \leqslant\left\Vert x-z\right\Vert +\left\Vert
z-y\right\Vert <t\left\Vert x-\xi\right\Vert +\delta<tR+t^{2}r<2tR=\varepsilon
.
\]
Hence $B(z,\delta)\subset C\cap B(x,\varepsilon)$. Thus in either case, there
is a ball of radius $\delta$ contained in $C\cap B(x,\varepsilon)$.
\end{proof}

\begin{lemma}
\label{0105a3}Let $\left\{  x_{1},\ldots,x_{N}\right\}  $ be a basis of
$\mathbb{R}^{N}$. Then $\tfrac{1}{N+1}%
%TCIMACRO{\tsum _{k=1}^{N}}%
%BeginExpansion
{\textstyle\sum_{k=1}^{N}}
%EndExpansion
x_{k}$ is in the interior of the convex hull of $\{0,x_{1},\ldots,x_{N}\}$.
\end{lemma}

\begin{proof}
We can unambiguously define a norm $\left\Vert \ \right\Vert ^{\prime}$ on
$\mathbb{R}^{N}$ by $\left\Vert \sum_{k=1}^{N}\lambda_{k}x_{k}\right\Vert
^{\prime}=\sum_{k=1}^{N}\left\vert \lambda_{k}\right\vert $. By the
equivalence of norms on finite-dimensional spaces, there exists $c>0$ such
that $\left\Vert \ \right\Vert ^{\prime}\leqslant c\left\Vert \ \right\Vert $,
where the latter norm is the original one on $X$. Let $x\equiv\sum_{k=1}%
^{n}\lambda_{k}x_{k}\in X$ satisfy%
\[
\ \left\Vert x-b_{n}\right\Vert <r\equiv\frac{1}{n(n+1)c}.
\]
Then%
\[
\sum_{k=1}^{n}\left\vert \lambda_{k}-\frac{1}{n+1}\right\vert =\left\Vert
x-b_{n}\right\Vert ^{\prime}<\frac{1}{n(n+1)},
\]
so for each $k$,%
\[
0<\frac{1}{n+1}-\frac{1}{n(n+1)}<\lambda_{k}<\frac{1}{n+1}+\frac{1}{n(n+1)}%
\]
and therefore%
\[
\sum_{k=1}^{n}\lambda_{k}<\frac{n}{n+1}+\frac{n}{n(n+1)}=1.
\]
Thus $x$ belongs to the convex hull of $\left\{  0,x_{1},\ldots,x_{n}\right\}
$. Hence $B\,(b_{n},r)$ in $\mathbb{R}^{N}$ is a subset of that convex hull.
\end{proof}

\begin{lemma}
\label{jul17l1}Let $\mathbf{S}$ be a measurable complemented set in
$\mathbb{R}^{N}$ with positive (possibly infinite) measure, such that $S^{1}$
is convex. Then the interior of $\ S^{1}$ is inhabited.
\end{lemma}

\begin{proof}
Translating $S^{1}$ if necessary, we may assume that\textbf{\ }$0\in S^{1}$.
By Lemma \ref{jul02l1}, $S^{1}$ contains a basis $\left\{  x_{1},\ldots
,x_{N}\right\}  $ of $\mathbb{R}^{N}$. By Lemma \ref{0105a3}, $b\equiv\frac
{1}{N+1}\sum_{n=1}^{N}x_{n}$\ belongs to the interior of the convex hull of
$\left\{  0,x_{1},\ldots,x_{N}\right\}  $ and hence to the interior of $S^{1}$.
\end{proof}

%

%TCIMACRO{\TeXButton{medskip}{\medskip}}%
%BeginExpansion
\medskip
%EndExpansion


Now we have our first major proof.%

%TCIMACRO{\TeXButton{medskip}{\medskip}}%
%BeginExpansion
\medskip
%EndExpansion


\begin{proof}
[Proof of Theorem 1]By Lemma \ref{jul17l1}, there exist $s_{0}\in S^{1}$ and
$r>0$ such that $B(s_{0},r)\subset S^{1}$. To begin with, assume
that\textbf{\ }$S^{1}$ is bounded. Then $\mathbf{S}$ is integrable, by Lemma
\ref{jul22l1}, and $\mu(\mathbf{S})\geqslant\mu(\mathbf{B}(s_{0},r))>0$. By
Lemma \ref{0105a2}, $\left(  S^{1}\right)  ^{\circ}$ is uniformly dense in
$S^{1}$. Given $\varepsilon>0$, pick $\delta>0$ such that if $x\in S^{1}$,
then $S^{1}\cap B(x,\varepsilon)$ contains a ball in $\mathbb{R}^{N}$ of
radius $\delta$. Applying \cite[Ch. 6, (6.7)]{BB}, construct a compact set $K$
such that $\mathbf{K}\equiv(K,-K)$ is integrable, $K\subset S^{1}$, and
$\mu(\mathbf{S}-\mathbf{K})<c$. Let $F$ be a finite $\varepsilon
$-approximation to $K$, and consider any $x\in S^{1}$. There exists $z\in
S^{1}$ such that $B(z,\delta)\subset S^{1}\cap B(x,\varepsilon)$. If
$\rho\left(  x,K\right)  >\varepsilon$, then $B(x,\varepsilon)\subset-K$, so
$B(z,\delta)\subset S^{1}-K$, $\mathbf{B}(z,\delta)<$ $\mathbf{S}-\mathbf{K}$,
and therefore $\mu\left(  \mathbf{S}-\mathbf{K}\right)  \geqslant
\mu(\mathbf{B}(z,\delta))$, a contradiction. Hence $\rho\left(  x,K\right)
<2\varepsilon$ and therefore $\rho(x,F)<3\varepsilon$. Thus $F$ is a finite
$3\varepsilon$-approximation to $S^{1}$. Since $\varepsilon$ is arbitrary, it
follows that $S^{1}$ is totally bounded and hence located.

Now consider the general case, in which $S^{1}$ need not be bounded. Given
$a\in\mathbb{R}^{N}$, choose%
\[
R>4\left\Vert a-s_{0}\right\Vert +2\left\Vert s_{0}\right\Vert
\]
such that $B(s_{0},r)\subset B(0,R)$, and let $\mathbf{A}=\mathbf{S}%
\wedge\mathbf{B}(0,R)$. Then $\mathbf{A}$ is integrable (by Lemma
\ref{jul21l1}), $\mathbf{B}(s,r)<\mathbf{A}$, and therefore $\mu
(\mathbf{A})\geqslant\mu(\mathbf{B}(s,r))>0$. Since $A^{1}$ is convex and
bounded, it follows from the first paragraph of this proof that $A^{1}$ is
located in $\mathbb{R}^{N}$. Consider any $x\in S^{1}$. Either $\left\Vert
x\right\Vert <R$, in which case $x\in A^{1}$ and $\left\Vert a-x\right\Vert
\geqslant\rho(a,A^{1})$, or else $\left\Vert x\right\Vert >R/2$. In the latter
case,%
\begin{align*}
\left\Vert a-x\right\Vert  &  \geqslant\left\Vert x\right\Vert -\left\Vert
a-s_{0}\right\Vert -\left\Vert s_{0}\right\Vert \\
&  >\frac{R}{2}-\left\Vert a-s_{0}\right\Vert -\left\Vert s_{0}\right\Vert \\
&  >\left\Vert a-s_{0}\right\Vert \\
&  \geqslant\rho(a,A^{1})\text{.}%
\end{align*}
Since $A^{1}\subset S^{1}$, it follows that $\rho(a,S^{1})$ exists and equals
$\rho(a,A^{1})$.
\end{proof}

\begin{corollary}
\label{jul17c1}Let $\mathbf{S}$ be a measurable complemented set with positive
(possibly infinite) measure in $\mathbb{R}^{N}$, such that $S^{1}$ is bounded
and convex. If $%
%TCIMACRO{\TeXButton{mathsim}{\mathord{\sim}}}%
%BeginExpansion
\mathord{\sim}%
%EndExpansion
S^{1}$ is inhabited, then it is located in $\mathbb{R}^{N}$.
\end{corollary}

\begin{proof}
By Lemma \ref{jul17l1}, $\left(  S^{1}\right)  ^{\circ}$ is inhabited, so,
$S^{1}$ being convex, $\left(  S^{1}\right)  ^{\circ}$ is dense in $S^{1}$
\cite[Propn 10]{dsb26}. Hence $-S^{1}=-\left(  S^{1}\right)  ^{\circ}$.
Applying Lemma 5.1.4 of \cite{BVtech} to the inhabited, open, convex set
$\left(  S^{1}\right)  ^{\circ}$, we see that $-S^{1}$ is dense in $%
%TCIMACRO{\TeXButton{mathsim}{\mathord{\sim}}}%
%BeginExpansion
\mathord{\sim}%
%EndExpansion
\left(  S^{1}\right)  ^{\circ}$. Since $-S^{1}\subset%
%TCIMACRO{\TeXButton{mathsim}{\mathord{\sim}}}%
%BeginExpansion
\mathord{\sim}%
%EndExpansion
S^{1}\subset%
%TCIMACRO{\TeXButton{mathsim}{\mathord{\sim}}}%
%BeginExpansion
\mathord{\sim}%
%EndExpansion
\left(  S^{1}\right)  ^{\circ}$, it follows that $-S^{1}$ is dense in $%
%TCIMACRO{\TeXButton{mathsim}{\mathord{\sim}}}%
%BeginExpansion
\mathord{\sim}%
%EndExpansion
S^{1}$. Referring to Theorem \ref{jul21t1}, we see that $S^{1}$ is a located
convex subset of $\mathbb{R}^{N}$ such that $-S^{1}$ is inhabited; whence, by
\cite[4.1.15]{BVtech}, the dense subset $-S^{1}$ of $%
%TCIMACRO{\TeXButton{mathsim}{\mathord{\sim}}}%
%BeginExpansion
\mathord{\sim}%
%EndExpansion
S^{1}$ is located in $\mathbb{R}^{N}$, from which the result follows.
\end{proof}

%

%TCIMACRO{\TeXButton{medskip}{\medskip}}%
%BeginExpansion
\medskip
%EndExpansion


\label{2207}In the case $N=1$ we can remove from Theorem \ref{jul21t1} the
hypothesis that $\mathbf{S}$ have positive (possibly infinite) measure. Here
is a simpler proof than the one given in \cite[Thm. 3]{Br88}.

\begin{proposition}
\label{jul21p1}Let $S$ be an inhabited, measurable, convex subset of
$\mathbb{R}$. Then $S^{1}$ is located in $\mathbb{R}$.
\end{proposition}

\begin{proof}
It will suffice to consider the case where $S^{1}$ is bounded, in which case
$\mathbf{S}$ is integrable: for we can then deal with the general case by a
simpler variant of the argument in the second paragraph of the proof of
Theorem \ref{jul21t1}. Pick $s\in S^{1}$. Given $\varepsilon>0$, we have
either $\mu(\mathbf{S})>0$, when Theorem \ref{jul21t1} applies, or else
$\mu(\mathbf{S})<\varepsilon/2$. In the latter case, if $x\in S^{1}$ and
$\left\vert x-s\right\vert >\varepsilon/2$, then since $S^{1}$, being convex,
contains all points of the interval between $s$ and $x$, so $\mu
(\mathbf{S})>\varepsilon/2$, a contradiction; whence $\left\vert
x-s\right\vert <\varepsilon$ for all $x\in S^{1}$ and therefore $\left\{
s\right\}  $ is an $\varepsilon$-approximation to $S^{1}$.
\end{proof}

%

%TCIMACRO{\TeXButton{medskip}{\medskip}}%
%BeginExpansion
\medskip
%EndExpansion


Here is a Brouwerian example\footnote{%
%TCIMACRO{\TeXButton{sf}{\normalfont\sf}}%
%BeginExpansion
\normalfont\sf
%EndExpansion
An improvement on the one on page 419 of \cite{Br88}.} showing that we cannot
drop the positive (possibly infinite) measure hypothesis in the case $N>1$
even if we know that $S^{1}$ is inhabited. Let $\left(  a_{n}\right)
_{n\geqslant1}$ be an increasing binary sequence with $a_{1}=0$. For each $n$
construct a compact, convex subset $S_{n}$ of $\mathbb{R}^{2}$ such that

\begin{itemize}
\item[--] if $a_{n}=0$, then $S_{n}=\left[  0,1\right]  \times\left\{
0\right\}  $;

\item[--] if $a_{n}=1-a_{n-1}$ and $n$ is odd, then $S_{k}=\left[  0,2\right]
\times\left\{  0\right\}  $ for all $k\geqslant n$;

\item[--] if $a_{n}=1-a_{n-1}$ and $n$ is even, then for all $k\geqslant n$,
$S_{n}$ is the closed triangle with vertices $\left(  0,0\right)  ,\left(
1,0\right)  ,\left(  1,2/n\right)  $.
\end{itemize}

%

%TCIMACRO{\TeXButton{noindent}{\noindent}}%
%BeginExpansion
\noindent
%EndExpansion
Then $\mathbf{S}_{n}\equiv\left(  S_{n},-S_{n}\right)  $ is integrable for
each $n$. Let $\mathbf{S}=%
%TCIMACRO{\tbigvee _{n\geqslant1}}%
%BeginExpansion
{\textstyle\bigvee_{n\geqslant1}}
%EndExpansion
\mathbf{S}_{n}$. Since $S_{n}\subset S_{n+1}$ for each $n$, $S^{1}$ is
inhabited and convex. To prove that $\mathbf{S}$ is integrable, for each $k$
let $\mathbf{A}_{k}$ be the integrable set $%
%TCIMACRO{\tbigvee _{n=1}^{k}}%
%BeginExpansion
{\textstyle\bigvee_{n=1}^{k}}
%EndExpansion
\mathbf{S}_{n}$; then $A_{1}^{1}\subset A_{2}^{1}\subset\cdots$. Moreover, if
$j\geqslant k$, then either $a_{j}=a_{k}$ and $\mu(\mathbf{A}_{j}%
)=\mu(\mathbf{A}_{k})$, or else $a_{k}=0$ and $a_{j}=1$. In the latter event,
$\mu(A_{k})=0$ and there exists $N$ such that $k<N\leqslant j$ and
$a_{N}=1-a_{N-1}$; if $N$ is odd, then $\mu(A_{j})=\mu(A_{N})=0$; if $N$ is
even, then $\mu(A_{j})=\mu(A_{N})=1/N$. Thus in all cases, $\left\vert
\mu(\mathbf{A}_{j})-\mu(\mathbf{A}_{k})\right\vert \leqslant1/k$. Hence
$\left(  \mu(\mathbf{A}_{k})\right)  _{k\geqslant1}$ is a Cauchy sequence and
therefore converges to a limit in $\mathbb{R}$. It follows from \cite[Ch. 6,
(3.8)]{BB} that $\mathbf{S}$ is integrable. On the other hand, if $S^{1}$ is
located, then either $\rho\left(  \left(  2,0\right)  ,S^{1}\right)  >0$ or
$\rho\left(  \left(  2,0\right)  ,S^{1}\right)  <1$. In the first case, if
$a_{m}=1-a_{m-1}$, then $m$ cannot be odd; in the second, if $a_{m}=1-a_{m-1}%
$, then $m$ cannot be even. We now see that if Theorem \ref{jul21t1} holds
without the hypothesis $\mu(\mathbf{S})>0$, then we can prove the following principle:

\begin{quote}
For each increasing binary sequence $\left(  a_{n}\right)  _{n\geqslant1}$
with $a_{1}=0$, either the least $m$, if there is one, with $a_{m}=1-a_{m-1}$
is even, or else the least such $m$ is odd.
\end{quote}

%

%TCIMACRO{\TeXButton{noindent}{\noindent}}%
%BeginExpansion
\noindent
%EndExpansion
This principle is easily shown to be equivalent to the well-known omniscience
principle\footnote{%
%TCIMACRO{\TeXButton{sf}{\normalfont\sf}}%
%BeginExpansion
\normalfont\sf
%EndExpansion
See page 9 of \cite{BVtech}.} \textsf{LLPO} and is therefore essentially nonconstructive.%

%TCIMACRO{\TeXButton{medskip}{\medskip}}%
%BeginExpansion
\medskip
%EndExpansion


We now turn to Theorem \ref{jul21t2}, which strengthens Corollary
\ref{jul17c1}. We shall need a simple variant of the constructive
least-upper-bound principle \cite[2.1.18]{BVtech}: namely, that if$\ A$ is an
inhabited subset of $\mathbb{R}$ such that for all $\beta\in\mathbb{R}$ either
$\inf A$ exists or else there exists $x\in A$ with $x<\beta$, then $\inf S$ exists.%

%TCIMACRO{\TeXButton{medskip}{\medskip}}%
%BeginExpansion
\medskip
%EndExpansion


\begin{proof}
[Proof of Theorem 2]Since $S^{1}\cup%
%TCIMACRO{\TeXButton{mathsim}{\mathord{\sim}}}%
%BeginExpansion
\mathord{\sim}%
%EndExpansion
S^{1}$, being a full set, is dense in $\mathbb{R}^{N}$, it will suffice to
prove that $\rho(x,%
%TCIMACRO{\TeXButton{mathsim}{\mathord{\sim}}}%
%BeginExpansion
\mathord{\sim}%
%EndExpansion
S^{1})$ exists for each $x$ in $S^{1}$. Given such $x$, from Lemma
\ref{jul21l1} we see that $\mathbf{S\wedge B}(x,r)$ is integrable for all
$r>0$. For each positive integer $n$ let $v_{n}=\mu(\mathbf{B}(0,n^{-1}%
))\ $and construct an increasing binary sequence $\left(  \lambda_{n}\right)
_{n\geqslant1}$ such that for each $n$,%
\begin{align*}
\lambda_{n}=0  &  \Rightarrow\mu(\mathbf{S}\wedge\mathbf{B}\left(
x,n^{-1}\right)  )<v_{n},\\
\lambda_{n}=1  &  \Rightarrow\mu(\mathbf{S}\wedge\mathbf{B}\left(
x,n^{-1}\right)  )>0.
\end{align*}
In view of Corollary \ref{jul17c1}, we may assume that $\lambda_{1}=0$. If
$\lambda_{n}=0$, then by \cite[Ch. 6, (3.6)]{BB},%
\[
\mu(\mathbf{B}\left(  x,n^{-1}\right)  -\mathbf{S})=v_{n}-\mu(\mathbf{S}%
\wedge\mathbf{B}\left(  x,n^{-1}\right)  )>0,
\]
so, by \cite[Ch. 6, (3.4)]{BB}, we can choose $x_{n}\in S^{0}\subset%
%TCIMACRO{\TeXButton{mathsim}{\mathord{\sim}}}%
%BeginExpansion
\mathord{\sim}%
%EndExpansion
S^{1}$ with $\left\Vert x_{n}-x\right\Vert <1/n$. If $\lambda_{n}=1$, we set
$x_{n}=x_{n-1}$. If $m\geqslant n$, then%
\[
\left\Vert x_{m}-x_{n}\right\Vert \leqslant\left\Vert x_{m}-x\right\Vert
+\left\Vert x_{n}-x\right\Vert <\frac{1}{m}+\frac{1}{n}\leqslant2/n;
\]
so $\left(  x_{n}\right)  _{n\geqslant1}$ is a Cauchy sequence and therefore
converges to a limit $x_{\infty}$ in the closure of $%
%TCIMACRO{\TeXButton{mathsim}{\mathord{\sim}}}%
%BeginExpansion
\mathord{\sim}%
%EndExpansion
S^{1}$. Consider any $\alpha,\beta$ with $0<\alpha<\beta$. Either $\left\Vert
x_{\infty}-x\right\Vert >\alpha$ or $\left\Vert x_{\infty}-x\right\Vert
<\beta$. In the first case, there exists $m$ such that $\left\Vert
x_{m}-x\right\Vert >1/m$; then $\lambda_{m}=1$, $\mu(\mathbf{S})\geqslant$
$\mu(\mathbf{S}\wedge\mathbf{B}\left(  x,m^{-1}\right)  )>0$, and by Corollary
\ref{jul17c1}, $%
%TCIMACRO{\TeXButton{mathsim}{\mathord{\sim}}}%
%BeginExpansion
\mathord{\sim}%
%EndExpansion
S^{1}$ is located. In the second case, there exists $n$ such that $\left\Vert
x_{n}-x\right\Vert <\beta$, where $x_{n}\in%
%TCIMACRO{\TeXButton{mathsim}{\mathord{\sim}}}%
%BeginExpansion
\mathord{\sim}%
%EndExpansion
S^{1}$. It follows from the remark immediately preceding this theorem that $%
%TCIMACRO{\TeXButton{mathsim}{\mathord{\sim}}}%
%BeginExpansion
\mathord{\sim}%
%EndExpansion
S^{1}$ is located.
\end{proof}

%

%TCIMACRO{\TeXButton{medskip}{\medskip}}%
%BeginExpansion
\medskip
%EndExpansion


In view of our earlier Brouwerian example, we see from Theorem \ref{jul21t2}
that if, in Theorem \ref{jul21t1}, we replace the hypothesis `$\mu
(\mathbf{S})>0$' by `$%
%TCIMACRO{\TeXButton{mathsim}{\mathord{\sim}}}%
%BeginExpansion
\mathord{\sim}%
%EndExpansion
S^{1}$ is inhabited', then we can prove that the latter set is located even if
we don't know whether $S^{1}$ itself is located.%

%TCIMACRO{\TeXButton{bigskip}{\bigskip}}%
%BeginExpansion
\bigskip
%EndExpansion


\begin{thebibliography}{9}                                                                                                %


\bibitem {Bishop}Errett Bishop, \emph{Foundations of Constructive Analysis},
McGraw-Hill, New York, 1967.

\bibitem {BB}E. Bishop and D.S. Bridges, \emph{Constructive Analysis},
Grundlehren der math. Wissenschaften \textbf{279}, Springer Verlag,
Heidelberg-Berlin-New York, 1985.

\bibitem {dsb26}D.S. Bridges,\ `Metric double complements of convex sets',
2025. https:// doi.org/10.48550/arXiv.2510.15123.

\bibitem {Br88}D.S. Bridges, `Locatedness, convexity, and Lebesgue
measurability', Quart. J. Math. Oxford (2) \textbf{39}, 411--421, 1988.

\bibitem {BVtech}D.S. Bridges, L.S. V\^{\i}\c{t}\u{a}, \emph{Techniques of
Constructive Analysis}, Universitext, Springer New York, 2006.

\bibitem {Handbook}D.S. Bridges, H. Ishihara, M.J. Rathjen, H. Schwichtenberg
(editors),\emph{\ Handbook of Constructive Mathematics}, Encyclopedia of
Mathematics and Its Applications \textbf{185}, Cambridge University Press, 2023.
\end{thebibliography}

%

%TCIMACRO{\TeXButton{bigskip}{\bigskip}}%
%BeginExpansion
\bigskip
%EndExpansion
%

%TCIMACRO{\TeXButton{bigskip}{\bigskip}}%
%BeginExpansion
\bigskip
%EndExpansion
%

%TCIMACRO{\TeXButton{bigskip}{\bigskip}}%
%BeginExpansion
\bigskip
%EndExpansion
%

%TCIMACRO{\TeXButton{noindent}{\noindent}}%
%BeginExpansion
\noindent
%EndExpansion
\textbf{{\small Author's address:}}{\small \ Department of Mathematics \&
Statistics, University of Canterbury, Christchurch 8140, New Zealand}%

%TCIMACRO{\TeXButton{noindent}{\noindent}}%
%BeginExpansion
\noindent
%EndExpansion
\textbf{{\small Author's email:} \ \thinspace}%
\texttt{{\small dugbridges@gmail.com}}

\bigskip

\vfill


\begin{flushright}
{\small \texttt{dsbridges 23 July 2026}}
\end{flushright}


\end{document}
